\section{Introduction}
\label{sec:intro}

Publicly supported co-creation hubs, enterprise hubs, and related collaborative spaces have become increasingly visible instruments of regional development policy. Local governments often view such hubs as a way to reduce interaction frictions among firms, researchers, entrepreneurs, and other regional actors by creating spaces and programs that facilitate meetings, information exchange, and collaboration. This policy logic is especially attractive in thin regional ecosystems, where the local pool of potentially complementary actors is limited and decentralized coordination may otherwise be weak. Yet the welfare case for hub policy is not obvious in such places. A hub may generate visible activity while still failing to create enough durable novelty to justify its cost, especially if the same local actors repeatedly interact within a narrow regional pool.

This tension sits naturally within the broader literature on agglomeration and regional coordination. A long tradition in urban and regional economics emphasizes that spatial concentration can raise productivity through matching, learning, and related external economies \citep{DurantonPuga2004,RosenthalStrange2004}. Face-to-face interaction has often been treated as a particularly important channel through which localized economic activity generates value \citep{StorperVenables2004}. Related work on local buzz and global pipelines similarly suggests that geographically proximate interaction can be productive, especially when local contact is complemented by connections to outside knowledge and opportunities \citep{BatheltMalmbergMaskell2004}. More recent evidence from co-working environments also indicates that proximate social interaction can indeed create knowledge spillovers and collaboration gains \citep{RocheOettlCatalini2024}. Taken together, these contributions provide a prima facie rationale for why a public hub might help a region.

At the same time, the same literature also implies that more local interaction is not always better. Proximity can support learning, but it can also generate lock-in when interaction becomes repetitive and the same actors repeatedly draw on the same local knowledge base \citep{Boschma2005,TorreRallet2005}. The regional innovation systems literature likewise stresses that local development depends not only on whether interaction occurs, but also on whether regional systems are able to renew themselves through access to external ideas, organizations, and opportunities \citep{CookeUrangaEtxebarria1997,AsheimLawtonSmithOughton2011,TodtlingTrippl2005}. This concern is especially relevant for place-based policy. Local interventions may correct coordination problems, but their effects are conditional and context-dependent rather than automatic \citep{KlineMoretti2014ARE,NeumarkSimpson2015}. In the case of public hubs, this raises a simple but underexplored question: when should a hub be evaluated as a stand-alone project, and when should it instead be evaluated against an active renewal strategy that already builds outside links?

This paper develops a stylized two-region, two-period framework for public co-creation hubs in thin regional ecosystems. Region $A$ is the target thin region in which a public hub may raise the local matching rate among potential collaborators. Region $B$ is an outside pool from which additional collaborators, ideas, and partnerships may arrive. It should not be interpreted as a fully modeled second regional economy. Rather, it is a reduced-form outside access margin through which cross-regional renewal can occur. The model distinguishes two policy margins. The first is \emph{local coordination}: a hub raises the rate at which potentially valuable local interactions in region $A$ are realized. The second is \emph{cross-regional renewal}: pipeline effort creates additional links between region $A$ and outside actors in region $B$.

A key modeling choice is to distinguish decentralized local support from social evaluation. The paper focuses on a \emph{dynamic coordination externality} that arises when current hub services and directly brokered outside links are contractible, but the future novelty created inside the regional system is not. Repeated local matching may exhaust novelty, while outside renewal can partially replenish it. The resulting wedge does not rely on an ad hoc short-run rule, but on a stylized appropriability problem in which future local novelty behaves like a diffuse regional asset. The setup is dynamic only in a minimal two-period sense: it isolates one intertemporal wedge rather than modeling cumulative lock-in in full generality.

The framework delivers three conditional implications. First, decentralized support and planner evaluation need not coincide. Given the maintained appropriability structure, decentralized actors underprovide renewal effort because they do not fully internalize its contribution to future local novelty, and hub support can be either too weak or too strong depending on whether renewal gains or redundancy losses dominate. Second, a hub is socially efficient as a stand-alone intervention only under limited conditions: its local coordination gain must be large enough to offset fixed cost and the dynamic redundancy generated by repeated local interaction. Third, pipeline policy has positive stand-alone value, but the relevant policy object is often a \emph{bundle} combining local coordination with cross-regional renewal. A bundle can dominate pipeline policy alone, but only under specific parameter conditions.

These results speak directly to current regional policy practice. In many settings, hubs are assessed using visible indicators such as participation counts, meetings, events, occupancy rates, or early project volume. The model suggests that such measures are informative about immediate coordination gains, but incomplete as welfare metrics. What matters for regional efficiency is not only whether a hub increases local interaction, but also whether novelty persists over time, whether outside renewal is adequately supplied, and whether the hub creates enough value relative to a pipeline-only alternative.

The paper's contribution is narrower than a general theory of regional development, but more pointed in policy terms. First, it redefines the relevant policy baseline: hub policy should often be evaluated not only against no intervention, but also against pipeline-only renewal. Second, it provides a minimal model that separates \emph{local coordination} from \emph{cross-regional renewal} as distinct policy margins. Third, it shows why visible local activity is an incomplete welfare metric when current coordination changes future novelty. In that sense, the paper speaks to the agglomeration, renewal, and place-based policy literatures by isolating a specific policy comparison they do not solve directly \citep{DurantonPuga2004,BatheltMalmbergMaskell2004,Boschma2005,KlineMoretti2014ARE,NeumarkSimpson2015}.

The model also generates empirically testable implications. It predicts that hub creation should raise initial interaction volume, but not necessarily preserve the share of new ties; that pipeline-only regions should exhibit more outside collaboration than no-policy regions; and that bundled regions should outperform pipeline-only regions when the outside pool is richer or regional absorptive capacity is stronger. These predictions point toward a dynamic and network-aware empirical evaluation of hub policy, one that tracks repeated ties, outside linkages, and participant composition rather than relying only on visible activity measures.

The remainder of the paper proceeds as follows. Section~\ref{sec:lit} reviews the related literature and situates the paper's contribution. Section~\ref{sec:model} presents the two-region model and the dynamic coordination externality. Section~\ref{sec:hub_efficiency} characterizes the conditions under which a hub is socially efficient as a stand-alone policy. Section~\ref{sec:short_run} compares decentralized support with planner evaluation. Section~\ref{sec:refresh} analyzes pipeline-alone policy, the bundled policy design, and the conditions under which the bundle is preferred. Section~\ref{sec:extensions} develops extensions on partial internalization and participant composition. Section~\ref{sec:policy} derives empirical predictions and discusses regional policy interpretation. Section~\ref{sec:conclusion} concludes.
